% ===============================
% SEZIONE 1: Parti di piano
% ===============================
\begin{question}[ID=1.1, subtitle={Rappresentazione di parti di piano — serie 1}, tags={parti di piano}, difficulty=facile]
Rappresenta graficamente le parti di piano definite dalle seguenti disequazioni o da sistemi di disequazioni:
\vspace{-10pt}
\begin{center}\begin{minipage}{0.7\linewidth}
\begin{enumerate}
\item $2 x - y > 1$
\item $\begin{cases} x + y - 2 = 0 \\ x \ge 1 \end{cases}$
\item $\begin{cases} 2x+y> 0 \\ y \ge 3x+1 \end{cases}$
\end{enumerate}
\end{minipage}\end{center}
\vspace{-20pt}
\begin{solution}
Si traccia la retta associata a ogni disequazione e si sceglie il semipiano corretto (continuo se $\le/\ge$, tratteggiato se $</>$).
\begin{itemize}
\item a) $2x-y>1\Rightarrow y<2x-1$: semipiano \emph{sotto} la retta $y=2x-1$ (esclusa).
\item b) $\{x+y-2=0,\ x\ge1\}$: la parte della retta $x+y=2$ con $x\ge1$, cioè la \emph{semiretta} a partire da $(1,1)$.
\item c) $\{2x+y>0,\ y\ge3x+1\}$: intersezione del semipiano sopra $y=-2x$ con quello sopra $y=3x+1$ (regione comune).
\end{itemize}
\end{solution}
\end{question}

\begin{question}[ID=1.2, subtitle={Rappresentazione di parti di piano — serie 2}, tags={parti di piano}, difficulty=facile]
Rappresenta graficamente le parti di piano definite dalle seguenti disequazioni o da sistemi di disequazioni:
\vspace{-10pt}
\begin{center}\begin{minipage}{0.7\linewidth}
\begin{enumerate}
\item $3 x - y < 1$
\item $\begin{cases} 2x + y - 2 = 0 \\ x \ge 1 \end{cases}$
\item $\begin{cases} 2x-y> 0 \\ y \le 3x+1 \end{cases}$
\end{enumerate}
\end{minipage}\end{center}
\vspace{-20pt}
\begin{solution}
\begin{itemize}
\item a) $3x-y<1\Rightarrow y>3x-1$: semipiano \emph{sopra} la retta $y=3x-1$ (esclusa).
\item b) $\{2x+y-2=0,\ x\ge1\}$: la parte della retta $2x+y=2$ con $x\ge1$, cioè la \emph{semiretta} a partire da $(1,0)$.
\item c) $\{2x-y>0,\ y\le3x+1\}$: intersezione del semipiano sotto $y=2x$ (cioè $y<2x$) con quello sotto/sulla retta $y=3x+1$.
\end{itemize}
\end{solution}
\end{question}


% ===============================
% SEZIONE 2: Problemi
% ===============================
\begin{question}[ID=2.1, subtitle={Punteggio di un test a risposta multipla}, tags={parti di piano}, difficulty=media]
Uno studente svolge un test a risposta multipla in cui il punteggio assegnato è pari a 1 punto per ogni risposta corretta, – 0,25 punti per ogni risposta errata e 0 punti per ogni risposta non data.
Il test è composto da 10 domande e, una volta completato, lo studente, confrontandosi con i compagni, si rende conto di aver sicuramente sbagliato almeno una domanda.
Rappresenta sul piano cartesiano ed elenca tutti i possibili esiti del test. \textit{(Suggerimento: poni come variabile $x$ il numero di domande corrette e $y$ il numero di domande errate)}
\begin{solution}
Con $x$ corrette e $y$ errate (interi $\ge0$), i vincoli sono: $x+y\le10$ (le restanti sono non date), $y\ge1$ (almeno una errata), $x\ge0$. La regione è il triangolo di vertici $(0,1)$, $(9,1)$, $(0,10)$, e gli esiti possibili sono i suoi \emph{punti a coordinate intere}.

Il punteggio associato a ciascun esito è $P=x-0{,}25\,y$. Ad esempio: $9$ corrette e $1$ errata $\to P=8{,}75$ (massimo possibile con almeno un errore); $0$ corrette e $10$ errate $\to P=-2{,}5$ (minimo). In generale, per ogni coppia $(x,y)$ con $x\ge0,\ y\ge1,\ x+y\le10$ si ottiene un esito ammissibile.
\end{solution}
\end{question}

\begin{question}[ID=2.2, subtitle={Composizione di un vassoio di dolci}, tags={parti di piano}, difficulty=media]
Una pasticceria prepara un vassoio di dolci composto da dolci al cioccolato e dolci alla frutta, per un totale di più di 8 dolci.
Si richiede che nel vassoio ci siano meno di 6 dolci alla frutta e che il numero dei dolci al cioccolato sia al massimo 5 in più rispetto a quelli alla frutta.
Rappresenta graficamente ed elenca tutti i possibili modi in cui può essere composto il vassoio di dolci.
\begin{solution}
Sia $x$ il numero di dolci al cioccolato e $y$ quelli alla frutta (interi $\ge0$). Vincoli: $x+y>8$ (cioè $x+y\ge9$), $y<6$ (cioè $y\le5$), $x\le y+5$. La condizione $9-y\le x\le y+5$ ha soluzioni solo se $9-y\le y+5$, cioè $y\ge2$. Quindi $y\in\{2,3,4,5\}$:
\begin{itemize}
\item $y=2$: $x=7$ $\to$ $(7,2)$;
\item $y=3$: $x=6,7,8$;
\item $y=4$: $x=5,6,7,8,9$;
\item $y=5$: $x=4,5,6,7,8,9,10$.
\end{itemize}
In totale $1+3+5+7=16$ composizioni possibili. Graficamente sono i punti interi della regione delimitata da $x+y\ge9$, $y\le5$, $x\le y+5$ (con $x,y\ge0$).
\end{solution}
\end{question}
